% Part: proof-theory
% Chapter: proof-search
% Section: completeness

\documentclass[../../../include/open-logic-section]{subfiles}

\begin{document}

\olfileid{pt}{ps}{cpl}

\olsection{تمامیت}

اکنون نشان می‌دهیم الگوریتم جست‌وجوی برهان خطاناپذیر است؛ یعنی اگر
متوقف شود یا هرگز پایان نیابد، برای دنبالهٔ آغازین
~$\Gamma \Sequent \Delta$ هیچ !!{proof}ی وجود ندارد. برای این کار
!!a{structure}~$\Struct M$ را چنان تعریف می‌کنیم که برای هر
$!A \in \Gamma$، $\Sat{M}{!A}$، و برای هر~$!A \in \Delta$،
$\Sat/{M}{!A}$. به بیان دیگر، همهٔ !!{formula}s موجود در~$\Gamma$ در
~$\Struct M$ صادق و همهٔ !!{formula}s موجود در~$\Delta$ کاذب‌اند؛ پس
$\Gamma \Sequent \Delta$ معتبر نیست. چون \Log{G3c} درست است،
$\Gamma \Sequent \Delta$ هیچ !!{proof}ی ندارد.

~‌$\Struct{M}$ را از یک جفت مجموعهٔ $\Theta, \Xi$ از !!{formula}s و
یک مجموعهٔ~$C$ از !!{constant}s تعریف می‌کنیم. $\Theta$، $\Xi$ و~$C$
را از یک \emph{شاخهٔ شکست} در اجرای ناموفق (متوقف‌شده یا پایان‌ناپذیر)
جست‌وجوی برهان به دست می‌آوریم. شاخهٔ شکست دنباله‌ای متناهی یا
نامتناهی از دنباله‌های $\Pi_n \Sequent \Lambda_n$ و مجموعه‌های
!!{constant}s~$C_n$ است، به‌گونه‌ای که $\Pi_0 \Sequent \Lambda_0$
همان دنبالهٔ پایانی $\Gamma \Sequent \Delta$ باشد و، برای هر~$n$ که
جملهٔ نهایی شاخهٔ متناهی نیست،
$\Pi_{n+1} \Sequent \Lambda_{n+1}$ مقدمه‌ای از
$\Pi_n \Sequent \Lambda_n$ باشد که در مرحلهٔ~$n+1$ تولید شده و
الگوریتم برای آن !!a{proof}ی نمی‌یابد. برای هر چنین~$n$ باید آشکارا
این مقدمه وجود داشته باشد، وگرنه الگوریتم !!a{proof}ی یافته بود.
$C_n$ را مجموعهٔ !!{constant}s متناظر با
~$\Pi_n \Sequent \Lambda_n$ می‌گیریم.

اگر الگوریتم در دنباله‌ای بالایی که تنها از !!{formula}s اتمی تشکیل
شده متوقف شود، شاخهٔ منتهی به آن دنباله یک شاخهٔ شکست است. اگر
الگوریتم هرگز پایان نیابد، پیوسته درخت‌های بزرگ‌تر تولید می‌کند. می‌توان
این درخت‌ها را مراحل رشد یک درخت نامتناهی دانست (درختی که پس از اجرای
بی‌نهایت مرحلهٔ الگوریتم به دست می‌آمد). این درخت نامتناهی، متناهی‌شاخه
است: هر گره تنها یک یا دو فرزند، یعنی دنباله‌های بی‌درنگ بالای خود،
دارد. بنا بر لم کونیگ، هر درخت نامتناهی متناهی‌شاخه یک شاخهٔ نامتناهی
دارد. اگر الگوریتم جست‌وجو تا ابد اجرا شود، چنین شاخه‌ای یک شاخهٔ شکست
خواهد بود.

اگر دنبالهٔ $\Pi_n \Sequent \Lambda_n$ و $C_n$ یک شاخهٔ شکست باشد،
می‌گذاریم $\Theta = \bigcup_n \Pi_n$ و
$\Xi = \bigcup_n \Lambda_n$ (با حذف اندیس‌ها و وقوع‌های تکراری
!!{formula}s)، و $C = \bigcup_n C_n$.

اکنون آماده‌ایم~$\Struct{M}$ را تعریف کنیم.
\begin{enumerate}
\item دامنهٔ $\Domain{M}$ مجموعهٔ ترم‌هایی است که می‌توان آن‌ها را
از~$C$ و !!{function}s موجود در $\Gamma \Sequent \Delta$، اگر چنین
نمادهایی وجود داشته باشند، و بدون هیچ !!{variable}s ساخت.
\item اگر $c \in C$، آنگاه $\Assign{c}{M} = c$.
\item اگر $f$ یک !!{function} $m$-موضعی در
$\Gamma \Sequent \Delta$ باشد و
$t_1$, \dots, $t_m \in \Domain{M}$، آنگاه
$\Assign{f}{M}(t_1, \dots, t_m) = f(t_1, \dots, t_m)$.
\item اگر $R$ یک !!{predicate} $m$-موضعی در
$\Gamma \Sequent \Delta$ باشد، آنگاه
$\Assign{R}{M} = \Setabs{\tuple{t_1, \dots, t_m} \in
\Domain{M}^m}{R(t_1, \dots, t_m) \in \Theta}$.
\end{enumerate}
~‌$\Struct{M}$ یک \emph{مدل ترمی} است، زیرا دامنه‌اش مجموعه‌ای از
ترم‌هاست و !!{constant}s و !!{function}s چنان تفسیر می‌شوند که
$\Value{t}{M} = t$ تضمین شود. !!{formula}s اتمی در~$\Theta$ برای
تعریف $\Assign{R}{M}$ چنان به کار می‌روند که
$\Sat{M}{R(t_1, \dots, t_m)}$ اگر و تنها اگر
$R(t_1, \dots, t_m) \in \Theta$ تضمین شود.

\begin{lem}\ollabel{lem:truth}
برای همهٔ $!A \in \Theta \cup \Xi$:
\begin{enumerate}
  \item اگر $!A \in \Theta$، آنگاه $\Sat{M}{!A}$.
  \item اگر $!A \in \Xi$، آنگاه $\Sat/{M}{!A}$.
\end{enumerate}
\end{lem}

\begin{proof}
  بر $\depth{!A}$ استقرا می‌کنیم.

نخست فرض کنید $!A$ اتمی است؛ یعنی یا $!A \ident \lfalse$ یا
$!A \ident R(t_1, \dots, t_m)$.

چون $\Pi_n \Sequent \Lambda_n$ شاخهٔ شکست است، هیچ $\Pi_n$ای
نمی‌تواند شامل $\lfalse$ باشد (وگرنه
$\Pi_n \Sequent \Lambda_n$ یک اصل بود). بنا بر تعریف~$\Struct{M}$،
$\Sat{M}{R(t_1, \dots, t_m)}$ اگر و تنها اگر
$R(t_1, \dots, t_m) \in \Theta$. این دو نکته با هم نشان می‌دهند اگر
$!A \in \Theta$، آنگاه $\Sat{M}{!A}$.

اگر $!A$ اتمی و $!A \in \Pi_n$ باشد، آنگاه
$!A \in \Pi_{n+1}$؛ و اگر $!A \in \Lambda_n$ باشد، آنگاه
$!A \in \Lambda_{n+1}$. (این امر از شیوهٔ تولید دنباله‌های جانشین در
الگوریتم جست‌وجوی برهان نتیجه می‌شود.) پس اگر
$!A \in \Theta \cap \Xi$، برای یک~$n$ داریم
$!A \in \Pi_n \cap \Lambda_n$. در این صورت
$\Pi_n \Sequent \Lambda_n$ اصل است، حال آنکه شاخهٔ شکست نمی‌تواند
اصلی داشته باشد. بنابراین بنا به ساخت، برای $!A$ اتمی
$!A \notin \Theta \cap \Xi$؛ پس اگر $!A \in \Xi$، آنگاه
$!A \notin \Theta$. بنا بر تعریف~$\Struct{M}$، از این امر نتیجه
می‌شود که اگر $!A \in \Xi$، آنگاه $\Sat/{M}{!A}$.

برای گام استقرا، فرض کنید (1) و (2) برای همهٔ !!{formula}s با
!!{depth} کمتر از~$!A$ برقرارند. (1) و (2) را برای~$!A$ با جداسازی
حالت‌ها اثبات می‌کنیم.

نخست توجه کنید که وقتی $!A^i$ در
$\Pi_n \Sequent \Lambda_n$ رخ دهد، در
$\Pi_{n+1} \Sequent \Lambda_{n+1}$ نیز رخ می‌دهد، مگر آنکه $i$
کوچک‌ترین اندیس در میان وقوع‌های غیراتمی
$\Pi_n \Sequent \Lambda_n$ باشد. چون اندیس‌های تازه همواره از اندیس
بیشینهٔ کنونی بزرگ‌ترند، سرانجام به
$\Pi_n \Sequent \Lambda_n$ای می‌رسیم که $!A^i$ در آن رخ می‌دهد و
$i$ کوچک‌ترین اندیس غیراتمی است. در مرحلهٔ $n+1$ الگوریتم~$!A^i$ را
کاهش می‌دهد. به بیان دیگر، اگر $!A \in \Theta$، برای یک $n$ داریم
$\Pi_n = \Pi_n', !A^i$ و $i$ کوچک‌ترین اندیس غیراتمی است؛ و اگر
$!A \in \Xi$، برای یک $n$ داریم
$\Lambda_n = \Lambda_n', !A^i$ و $i$ کوچک‌ترین اندیس غیراتمی است.

\begin{enumerate}
  \item فرض کنید $!A \in \Theta$ و
  $!A \ident !B \land !C$. برای یک $n$ داریم
  $\Pi_n = \Pi_n', (!B \land !C)^i$.
  $\Pi_{n+1} \Sequent \Lambda_{n+1}$ مقدمهٔ متناظر
  ~\LeftR{\land} است؛ یعنی
  $\Pi_{n+1} = \Pi_n', !B^{k+1}, !C^{k+2}$. در نتیجه
  $!B, !C \in \Theta$. بنا بر فرض استقرا، $\Sat{M}{!B}$ و
  $\Sat{M}{!C}$. بنا بر تعریف $\Sat{M}{}$، داریم
  $\Sat{M}{!B \land !C}$.

  \item فرض کنید $!A \in \Xi$ و $!A \ident !B \land !C$. برای
  یک $n$ داریم $\Lambda_n = \Lambda_n', !B \land !C$.
  $\Pi_{n+1} \Sequent \Lambda_{n+1}$ یکی از دو مقدمهٔ
  ~\RightR{\land} با نتیجهٔ
  $\Pi_n \Sequent \Lambda'_n, !B \land !C$ است؛ یعنی
  $\Lambda_{n+1} = \Lambda_n', !B^{k+1}$ یا
  $\Lambda_{n+1} = \Lambda_n', !C^{k+1}$. در نتیجه یا
  $!B \in \Xi$ یا $!C \in \Xi$. بنا بر فرض استقرا، یا
  $\Sat/{M}{!B}$ یا $\Sat/{M}{!C}$. بنا بر تعریف $\Sat{M}{}$، داریم
  $\Sat/{M}{!B \land !C}$.

  \item حالت‌های $!A \ident \lnot !B$، $!A \ident !B \lor !C$ و
  $!A \ident !B \lif !C$ مشابه‌اند و به‌عنوان تمرین واگذار می‌شوند.

  \item فرض کنید $!A \in \Theta$ و
  $!A \ident \lexists[x][!B(x)]$. برای یک $n$ داریم
  $\Pi_n = \Pi_n', \lexists[x][!B(x)]^i$.
  $\Pi_{n+1} \Sequent \Lambda_{n+1}$ مقدمهٔ متناظر
  ~\LeftR{\lexists} است؛ یعنی
  $\Pi_{n+1} = \Pi_n', !B(c)^{k+1}$ برای یک~$c$. در نتیجه
  $!B(c) \in \Theta$ و~$c \in C$. بنا بر فرض استقرا
  $\Sat{M}{!B(c)}$. بنا بر تعریف $\Sat{M}{}$، داریم
  $\Sat{M}{\lexists[x][!B(x)]}$.

  \item فرض کنید $!A \in \Xi$ و
  $!A \ident \lexists[x][!B(x)]$. برای یک $n$ داریم
  $\Lambda_n = \Lambda_n', \lexists[x][!B(x)]^i$. هر بار که
  $\lexists[x][!B(x)]$ کاهش می‌یابد، با اندیسی تازه در سمت راست مقدمه
  باقی می‌ماند؛ بنابراین اگر در شاخهٔ شکست در سمت راست رخ دهد،
  بی‌نهایت بار در آن سمت کاهش می‌یابد. هر ترم
  $t \in \Domain{M}$ سرانجام در شاخهٔ شکست «نخستین ترمی خواهد بود که
  تنها شامل ثابت‌های~$C_n$ است و پیش‌تر در کاهش راست
  $\lexists[x][!B(x)]$ به کار نرفته است». پس برای هر
  $t \in \Domain{M}$، برای یک~$n$ داریم $!B(t) \in \Lambda_n$ و در
  نتیجه $!B(t) \in \Xi$. بنا بر فرض استقرا، برای همهٔ
  $t \in \Domain{M}$ داریم $\Sat/{M}{!B(t)}$. بنا بر تعریف
  $\Sat{M}{}$، نتیجه می‌شود
  $\Sat/{M}{\lexists[x][!B(x)]}$.

  \item حالت‌های $!A \ident \lforall[x][!B(x)]$ مشابه‌اند و
  به‌عنوان تمرین واگذار می‌شوند.
\end{enumerate}
\end{proof}

\begin{cor}\ollabel{cor:complete}
الگوریتم جست‌وجوی برهان برای \Log{G3c} کامل است: اگر متوقف شود یا
هرگز پایان نیابد، دنبالهٔ آغازین~$\Gamma \Sequent \Delta$ معتبر نیست
و بنابراین هیچ !!{proof}ی ندارد.
\end{cor}

\begin{proof}
  اگر الگوریتم متوقف شود یا هرگز پایان نیابد، شاخهٔ شکستی وجود دارد که
  می‌توان $\Struct{M}$ را از آن تعریف کرد. بنا بر \olref{lem:truth}،
  برای همهٔ $!A \in \Gamma$ داریم $\Sat{M}{!A}$ و برای همهٔ
  $!A \in \Delta$ داریم $\Sat/{M}{!A}$. (به یاد آورید
  $\Pi_0 = \Gamma$ و $\Lambda_0 = \Delta$؛ پس
  $\Gamma \subseteq \Theta$ و $\Delta \subseteq \Xi$.) بنابراین
  $\Gamma \Sequent \Delta$ معتبر نیست. بنا بر درستی~\Log{G3c}،
  $\Gamma \Sequent \Delta$ هیچ !!{proof}ی ندارد.
\end{proof}

\begin{cor}
~\Log{G3c} کامل است: اگر $\Gamma \Sequent \Delta$ معتبر باشد،
$\Log{G3c} \Proves \Gamma \Sequent \Delta$.
\end{cor}

\begin{proof}
  عکس نقیض را اثبات می‌کنیم. فرض کنید
  $\Log{G3c} \Proves/ \Gamma \Sequent \Delta$. چون هیچ !!{proof}ی
  برای یافتن وجود ندارد، الگوریتم جست‌وجوی برهان باید متوقف شود یا
  هرگز پایان نیابد. بنا بر \olref{cor:complete}،
  $\Gamma \Sequent \Delta$ معتبر نیست.
\end{proof}


\end{document}
